% slides.tex — companion talk deck for "The Limitations of Generative AI"
% Build with: latexmk -pdf slides.tex
%
% The deck follows the paper's structure and holds its discipline: every
% number on these slides comes from the paper's verified bibliography, and
% claims are marked observed / explained / narrated where the distinction
% carries weight. Attributions are inline author-year text; the paper is
% the citation of record.

\documentclass[11pt]{beamer}
\usetheme{metropolis}
\usepackage{booktabs}
\usepackage{appendixnumberbeamer}
\usepackage{pgfpages}
% slides-notes.pdf = audience slide + speaker notes side by side
\ifdefined\WITHNOTES
\setbeameroption{show notes on second screen=right}
\fi

\title{The Limitations of Generative AI}
\subtitle{What these systems cannot (yet) do, and why}
\author{Nadia Yvette Chambers}
\date{\today\\[0.3em]{\small \url{https://doi.org/10.5281/zenodo.22294860}}}

\begin{document}

\maketitle

\begin{frame}{The problem with how AI is explained}
  \begin{itemize}
    \item Public understanding is built on \alert{demonstrations of success}:
          successes are easy to see, the structure of failures is not.
    \item Someone who has only seen a system succeed cannot tell a reliable
          tool from an unreliable one, judge announced fixes, or assess
          breathless claims in either direction.
    \item This talk supplies the missing half of the picture, under one
          discipline: every claim is marked as \textbf{observed} (measured),
          \textbf{explained} (well-argued mechanism), or
          \textbf{narrated} (repeated until it sounds established).
    \item The bibliography was not assembled from memory: all 51 entries were
          verified against DBLP, the arXiv API, and Crossref by a script, and
          open-access copies of 29 cited papers are archived with pinned
          hashes. Two claims that failed verification were dropped, not cited.
  \end{itemize}
  \note{Open cold with the asymmetry: everyone in the room has seen demos;
    almost nobody has been shown a taxonomy of failures. The three labels
    --- observed, explained, narrated --- are the talk's contract, so spend
    real time here (~2 min). Flag the two dropped claims as proof the
    discipline bites: the bibliography survived verification, two anecdotes
    did not. If someone asks which claims were dropped: an NSA interception
    anecdote with no stable sourcing, and a mis-dated surveillance item;
    both recorded in VERIFICATION.md.}
\end{frame}

\begin{frame}{Three families of failure --- not one laundry list}
  \begin{table}
    \centering\footnotesize
    \begin{tabular}{@{}p{0.20\textwidth}p{0.34\textwidth}p{0.38\textwidth}@{}}
      \toprule
      \textbf{Family} & \textbf{Signature} & \textbf{Matching remedy} \\
      \midrule
      Statistical & errors average out over samples & better data, better measurement \\
      Dynamical   & errors circulate through feedback & separate learning from acting \\
      Normative   & the violated rule cannot even be represented & inspectable structure outside the weights \\
      \bottomrule
    \end{tabular}
  \end{table}
  \begin{alertblock}{The consequential category error}
    Offering a statistical remedy for a dynamical or normative failure ---
    the most common shape of announced fixes and AI policy debates alike.
  \end{alertblock}
  \note{This table is the spine of the talk --- everything after it fills
    in one row. Walk left to right: family, signature, matching remedy.
    Then land hard on the alert block: the category error is the thing to
    carry home. Concrete hook: ``more data'' offered as the fix for a
    privacy violation is exactly this error --- a statistical remedy for a
    normative failure. Promise the audience they can diagnose announced
    fixes with one question: which family does the failure belong to?
    (~2 min.)}
\end{frame}

\section{Statistical failures}

\begin{frame}{Confabulation is the training objective working as designed}
  \begin{itemize}
    \item ``Hallucination'' suggests a perception problem; the mechanism is
          the objective itself rewarding \alert{plausible-sounding continuation}.
    \item Observed: scaling \emph{dilutes} but does not remove it; retrieval
          patches the symptom on popular topics while the long tail stays
          exposed (Mallen et al., 2023).
    \item The state is \emph{malleable}: a masked-word test that changes one
          word of the prompt flips the answer it gives (Pride \& Prejudice
          marriage dates, .org/.com) --- what is ``known'' versus what is
          reported come apart.
  \end{itemize}
  \note{Set up the reframe first: ``hallucination'' implies a broken
    perception module, but the mechanism is the objective rewarding
    plausible continuation --- confabulation is the system working as
    designed. Mallen et al.: retrieval patches the popular tail; the long
    tail stays exposed. The malleability result is the most vivid datum ---
    changing ONE word flips the answer, so ``what is known'' and ``what is
    reported'' come apart. Watch for the heckle ``won't scale fix it?'':
    answer with the observed dilution-not-removal line, and note that
    retrieval moves rather than removes the exposure.}
\end{frame}

\begin{frame}{Context and calibration}
  \begin{block}{Context is a budget, not a capability}
    \begin{itemize}
      \item Effective vs.\ advertised length: measured gap (Hsieh et al.,
            2024, RULER); ``lost in the middle'': recall is high at the
            edges and collapses in between (Liu et al., 2024).
      \item Retrieval augmentation inherits the problem: whatever is
            retrieved poorly is silently absent.
    \end{itemize}
  \end{block}
  \begin{block}{Calibration}
    \begin{itemize}
      \item A system that is right 80\% of the time should say so;
            badly calibrated systems sound equally sure either way.
      \item Masked-word experiments show the disconnect between held
            knowledge and reported confidence directly.
    \end{itemize}
  \end{block}
  \note{Two distinct failures share a slide --- keep them separate. Context:
    RULER measured the gap between effective and advertised length; Liu's
    ``lost in the middle'' is the positional curve --- edges good, middle
    collapsed. The RAG point lands well with technical audiences: whatever
    retrieval misses is silently absent, and the user cannot see the hole.
    Calibration: right 80\% of the time should SOUND like 80\%; bad
    calibration means equal confidence on right and wrong answers. Bridge
    to next slide: ``and if you think the tests we measure this with are
    trustworthy --- they aren't.''}
\end{frame}

\begin{frame}{Even the measurement layer is compromised}
  \begin{itemize}
    \item \textbf{Contamination}: test questions leak into training data ---
          examinations become memory tests.
    \item Fresh GSM1k problems matching the GSM8k distribution: drops of up
          to \alert{8 percentage points} in several model families ---
          though, honestly, the frontier systems showed little of it
          (Zhang et al., 2024).
    \item So a headline benchmark number can conflate reasoning with
          memorization, and you cannot tell which from the score alone.
  \end{itemize}
  \note{Contamination first: leaked test questions turn exams into memory
    tests. GSM1k is the headline --- fresh problems, same distribution, up
    to 8 points down. DELIVER THE HONESTY CLAUSE: frontier models showed
    little drop. This is where the talk earns trust --- report the result
    that weakens the dramatic reading. The conclusion is subtle: you can't
    tell reasoning from memorization from the score alone, so the number
    itself is the unreliable artifact. If challenged on GSM1k's own
    validity: it matches GSM8k's distribution by construction; the paper
    documents the matching procedure.}
\end{frame}

\section{Dynamical failures}

\begin{frame}{Memory: the read side of state}
  \begin{itemize}
    \item A fixed context window meets unbounded interaction: long sessions
          degrade retrieval of their own beginnings --- forgetting what was
          said \emph{in this conversation}.
    \item This is architectural, not a data problem: no amount of better
          training data changes the shape.
    \item Remedies that actually match: explicit external state ---
          scratchpads, retrieval over interaction history, persistent stores
          --- with write policies that are themselves inspectable.
  \end{itemize}
  \note{Transition into the dynamical family: ``from errors that average
    out to errors that circulate.'' Frame: a fixed window meets unbounded
    interaction --- the session forgets its own BEGINNINGS. Emphasize:
    architectural, not a data problem --- no training corpus changes the
    shape of the window. The remedy side matters for the Q\&A later:
    external state with inspectable write policies. This is also where the
    audience's lived experience (``I have to repeat myself by turn forty'')
    confirms the mechanism --- invite that recognition.}
\end{frame}

\begin{frame}{Learning in the loop: when errors circulate}
  \begin{itemize}
    \item A system that learns while it acts turns isolated errors into
          \alert{oscillations}: feedback loops, reward hacking, objective
          gaming --- and \emph{model collapse} when models train on their own
          outputs and the tail of the distribution decays.
    \item The remedy is \textbf{architectural separation}: learning on one
          timescale, acting on another --- two-process designs, frozen
          weights with an external controller.
    \item Statistical remedies (more data, better prompts) do not touch
          this family at all.
  \end{itemize}
  \note{The loop slide: learning while acting turns isolated errors into
    oscillations. Name the phenomena: feedback loops, reward hacking,
    objective gaming; model collapse when models train on their own
    outputs and the tail decays. The remedy is architectural separation
    --- learning on one timescale, acting on another (two-process designs,
    frozen weights with an external controller). Close the section by
    calling back to the table: statistical remedies do not touch this
    family AT ALL. That sentence is the section's takeaway.}
\end{frame}

\section{Normative failures}

\begin{frame}{Some obligations are rules, not probabilities}
  \begin{itemize}
    \item Privacy (contextual integrity, Nissenbaum 2009) is a system of
          norms: who may transmit what about whom, to whom, in which context.
    \item Formalized, it becomes a deontic logic with temporal operators
          (Barth et al., 2006) --- requiring recipient beliefs, role and
          group structure, retention over time, leak probabilities.
    \item Read the machinery off the formalization: \alert{multi-type,
          multi-agent, multimodal logic with temporality and probabilism}.
    \item Nothing in a corpus gradient can represent ``this transmission
          violates this norm in this context.''
  \end{itemize}
  \note{The pivot slide --- slow down. Privacy via Nissenbaum's contextual
    integrity: who may transmit what about whom, to whom, in which
    context. Then the formal move: it becomes a deontic logic with
    temporal operators (Barth et al.). Read the machinery requirement off
    the formalization: multi-type, multi-agent, multimodal, temporal,
    probabilistic. The closing sentence is the thesis of the whole
    normative section: NOTHING in a corpus gradient can represent ``this
    transmission violates this norm in this context.'' Expect pushback
    ``models do refuse'': defer it explicitly to the next slide.}
\end{frame}

\begin{frame}{Instances are not coherence}
  \begin{itemize}
    \item A fine-tuned model sometimes refuses to repeat someone's address:
          unreliably, with no audit trail, no consistency across time, and
          no correction short of retraining.
    \item Coherence --- norm consistency, auditability, reversibility ---
          requires \textbf{inspectable structure outside the trained weights}.
    \item The semiotic layer matters here: an \emph{indexical} link (this
          photo \emph{is} this person) is a doxxing vector; an iconic
          rendering is representational harm; a symbolic misuse violates an
          identifiable norm.
  \end{itemize}
  \begin{block}{Deployment reality}
    Narrow hybrids (guardrails, constrained decoding) are commercial.
    Flagship neurosymbolic systems (AlphaGeometry, AlphaProof) are research
    stage. The full governance-demanding profile is deployed
    \alert{nowhere} --- the vacuum is the finding.
  \end{block}
  \note{Answer the deferred objection directly: yes, fine-tuned models
    sometimes refuse --- unreliably, no audit trail, no consistency, no
    correction short of retraining. Instances are not coherence. The
    semiotic aside is a gem if the audience is philosophical: an indexical
    link (this photo IS this person) is a doxxing vector; an iconic
    rendering is representational harm; a symbolic misuse violates an
    identifiable norm. Land the block: narrow hybrids are commercial,
    flagship neurosymbolic is research-stage, the full profile is deployed
    NOWHERE --- the vacuum is the finding.}
\end{frame}

\section{Alternatives in the design space}

\begin{frame}{Diffusion language models change the trade-offs}
  \begin{itemize}
    \item Generate \alert{all positions in parallel}, in any order; every
          token is conditioned on every other (bidirectional context).
    \item Observed: an 8B diffusion LM is competitive with an equivalently
          sized autoregressive model (LLaDA vs.\ LLaMA 3 8B).
    \item Caveats: unlearning pathological associations is harder when all
          tokens are mutually conditioned; latency wins come at scale.
    \item Scope: helps \emph{statistical} limits (all-position access for
          context, parallel generation) --- not the dynamical or normative
          families.
  \end{itemize}
  \note{Guard the scope of this slide or it will be misheard as ``diffusion
    fixes everything.'' What diffusion changes: all positions in parallel,
    bidirectional conditioning; LLaDA 8B is competitive with LLaMA 3 8B.
    Honest caveats: unlearning is HARDER when all tokens are mutually
    conditioned; latency wins come at scale. Then the scope line --- it
    helps the statistical family only. Connect to their intuitions:
    parallel generation is why this matters for latency and multimodality,
    but it is still a distribution learner; no deontic structure appears.
    This slide exists so the alternatives section doesn't look
    unexplored.}
\end{frame}

\section{Delivery, incentives, and misrepresentation}

\begin{frame}{The shape of the machine shapes the science}
  \begin{itemize}
    \item Weights withheld, service only: transparency scores are lowest
          exactly where a service model hides things --- training data,
          compute, hardware (Bommasani et al., 2023, FMTI).
    \item Hosted models are updated or retired without documentation, and
          published findings silently break (Pozzobon et al., 2023).
    \item Access to scale was concentrating before the boom (Ahmed,
          2020); ML-based science has a reproducibility crisis of its own
          (Kapoor \& Narayanan, 2023). Google's 2023 memo said the quiet
          part: no moat, and neither does OpenAI.
  \end{itemize}
  \note{Section pivot: from what the systems cannot do to how they are
    sold and hosted. FMTI: transparency is lowest exactly where the
    service model hides things (data, compute, hardware). Pozzobon: hosted
    models updated/retired without documentation --- published findings
    silently break, which is a reproducibility problem for the FIELD.
    Ahmed: concentration of access predates the boom; Kapoor \& Narayanan:
    ML-based science has its own leakage-and-reproducibility crisis. The
    memo line gets a dry delivery: ``Google's own researchers wrote the
    moat down: none --- for them or OpenAI.''}
\end{frame}

\begin{frame}{Why states now treat hosted AI as a dependency question}
  \begin{itemize}
    \item A hosted service runs on someone else's infrastructure and
          jurisdiction, and the history keeps arriving: Merkel's phone
          (2013), Crypto AG secretly CIA-owned (2020), Danish cables used
          against European politicians (2021), a nine-year campaign of
          surveillance and intimidation against ICC staff (2024),
          ECHELON, and \emph{Schrems II} (2020).
    \item The lesson generalizes: exposure is a property of
          \alert{who operates the infrastructure}, not of who holds which
          alliance today.
    \item Hence \emph{digitale Souver\"anit\"at}: procurement moves
          (Schleswig-Holstein, Denmark) and the Data Act's switching and
          portability obligations.
  \end{itemize}
  \note{Keep this factual and unhurried --- the incident trail does the
    work. Merkel's phone (2013), Crypto AG secretly CIA-owned (2020),
    Danish cables used against European politicians (2021), the nine-year
    surveillance-and-intimidation campaign against ICC staff (2024),
    ECHELON, Schrems II. Then the generalization --- the talk's key
    sovereignty sentence: exposure is a property of WHO OPERATES the
    infrastructure, not of who holds which alliance today. (The ICC item
    is deliberately multi-actor: it is not a single-adversary narrative.)
    Land on digitale Souver\"anit\"at: Schleswig-Holstein and Denmark
    procurement moves, the Data Act's switching/portability obligations.
    ~2 min.}
\end{frame}

\begin{frame}{The everyday side: what does the customer possess?}
  \begin{itemize}
    \item A self-contained product keeps working when its maker does not;
          a service that shuts down --- acquisition, repricing, lost
          interest, bankruptcy --- takes its function with it.
    \item Not distinctive to AI: the standard lifecycle risk of hosting ---
          which is exactly what the sovereignty programs legislate against.
    \item Self-hosting and open weights exist in the design space (the
          Mowgli stack is a small, deliberately modest existence proof);
          the question is why the dominant presentation is the one where
          the customer owns least.
  \end{itemize}
  \note{The everyday counterweight to the geopolitical slide: what does
    the customer POSSESS? A self-contained product outlives its maker; a
    service shutdown (acquisition, repricing, lost interest, bankruptcy)
    takes the function with it. Stress: not distinctive to AI --- the
    standard lifecycle risk of hosting, exactly what sovereignty programs
    legislate against. Name Mowgli as a deliberately modest existence
    proof, not a competitor. Close with the pointed question --- why is
    the dominant presentation the one where the customer owns least? Hold
    the answer for the incentives slides.}
\end{frame}

\begin{frame}{Misrepresentation needs no dishonesty}
  \begin{itemize}
    \item Each link is individually reasonable; the chain is not:
          vendors fund demonstrations of success, the visible scoreboard is
          run with the firms it ranks, and the products are tuned for
          approval.
    \item \textbf{Sycophancy} (Sharma et al., 2023): agreement with the
          user's stated belief across every assistant studied, traced to
          human preference data and amplified by preference-model
          optimization --- apparent agreement is not evidence of correctness.
    \item \textbf{Leaderboard Illusion} (Singh et al., 2025): up to
          \alert{27 private variants} tested in parallel by select
          providers; \alert{205 of 243} public models silently deprecated;
          restricted tuning worth up to \alert{112\%} on a downstream
          benchmark.
  \end{itemize}
  \note{The section's punchline --- and the sentence that keeps the paper
    defensible: misrepresentation needs no dishonesty. Each link is
    individually reasonable; the chain is not. Sycophancy (Sharma):
    agreement with stated belief across EVERY assistant studied, traced
    to preference data and amplified by preference-model optimization ---
    so apparent agreement is not evidence of correctness. Leaderboard
    Illusion (Singh) has the numbers to say slowly: up to 27 private
    variants tested in parallel by select providers; 205 of 243 public
    models silently deprecated; restricted tuning worth up to 112\% on a
    downstream benchmark. Let the numbers sit.}
\end{frame}

\begin{frame}{Honesty about the boundary}
  \begin{block}{Measured}
    The capability claims are shaped by the same incentives that selected
    the architecture: this half is evidenced.
  \end{block}
  \begin{block}{Open, falsifiable}
    Whether those incentives \emph{reduce} total progress or merely
    \emph{redistribute} it toward the measurable and saleable.
    \emph{Weakening conditions}: independent instruments anchoring public
    perception; replication failure of the asymmetries; their absence in
    deployed systems.
  \end{block}
  \note{THE boundary slide --- deliver it with total precision, this is
    where speakers overreach. Measured: capability claims are shaped by
    the same incentives that selected the architecture. Open and
    falsifiable: whether those incentives REDUCE total progress or merely
    REDISTRIBUTE it toward the measurable and the saleable. Read the
    weakening conditions aloud --- independent instruments anchoring
    perception, replication failure of the asymmetries, their absence in
    deployed systems. Explicitly say: ``this is the claim I would most
    like to be wrong about, and here is how you would show it.'' That
    sentence is the talk's credibility in miniature.}
\end{frame}

\section{Open questions}

\begin{frame}{What remains open}
  \begin{itemize}
    \item Is confabulation reducible to acceptable levels under an objective
          that rewards plausibility? (\emph{unresolved})
    \item Do effective and advertised context lengths converge?
          (\emph{persistent across generations})
    \item Is reliable compositionality reachable by training alone?
          (\emph{evidence points to structural support})
    \item Do the incentives slow progress, or redistribute it?
          (\emph{falsifiable; conditions stated})
  \end{itemize}
  \begin{block}{The discipline}
    In-principle limits and in-practice shortfalls are different claims,
    and conflating them is how both hype and dismissal manufacture
    their certainty.
  \end{block}
  \note{Recap the open problems with their status tags --- unresolved,
    persistent across generations, structural support pointed to,
    falsifiable. Do NOT overclaim any of them; the status tags are the
    point. Then the closing block, which is the second-most-quotable
    line of the talk: in-principle limits and in-practice shortfalls are
    DIFFERENT claims, and conflating them is how both hype and dismissal
    manufacture their certainty. Pause after it. This is the bridge into
    takeaways.}
\end{frame}

\section{Takeaways}

\begin{frame}{Takeaways}
  \begin{enumerate}
    \item Ask \textbf{which family} of failure you are looking at --- the
          remedy only exists in the matching family.
    \item Distrust single-number evaluations: fresh test sets, independent
          instruments, style-sensitivity checks.
    \item For deployed decisions, ask what you \textbf{possess}: weights,
          audit rights, switching paths --- or only a subscription.
    \item Coherence (privacy, audit, consistency) demands inspectable
          structure; a probability distribution cannot break a rule it
          cannot represent.
  \end{enumerate}
  \note{Four takeaways, one breath each: (1) ask which FAMILY of failure
    --- remedy exists only in the matching family; (2) distrust
    single-number evaluations --- fresh test sets, independent
    instruments, style-sensitivity checks; (3) for deployed decisions ask
    what you POSSESS --- weights, audit rights, switching paths, or only
    a subscription; (4) coherence demands inspectable structure --- a
    probability distribution cannot break a rule it cannot represent.
    If time is short, cut nothing here --- cut the diffusion slide and
    compress open questions instead. These four are what the audience
    should be able to reproduce tomorrow.}
\end{frame}

\begin{frame}
  \centering
  \vspace{1.5em}
  {\Large The Limitations of Generative AI}\\[0.8em]
  \begin{itemize}
    \item Paper (with verified bibliography and corrections log):
          \url{https://nadiayvette.frama.io/gen-ai-limits/main.pdf}
    \item Source on all mirrors; framagit is canonical:
          \url{https://framagit.org/NadiaYvette/gen-ai-limits}
  \end{itemize}
  \vspace{1em}
  Thank you.
  \note{Leave this up during questions --- the URLs are the call to
    action. Expected Q\&A clusters and one-line deflections: (1) ``So AI
    is useless?'' --- No; the claim is that failure structure is knowable
    and remedies are family-specific. (2) ``Isn't neurosymbolic too
    hard/slow?'' --- deployed-nowhere is the finding, not an argument
    that it must stay there. (3) ``Isn't this just anti-industry?'' ---
    the measured/explained/narrated labels apply to the critique too;
    the falsifiability slide states the weakening conditions. (4)
    Timing-based: offer the bibliography --- every number on these
    slides traces to a hash-pinned PDF.}
\end{frame}

\end{document}
